\cotefile{29}
\scanpage{119}
% archive sheet number: 118 (uncertain; alternative 119)

\handwritten{2}. Supposons donné une situation $(Y,V,H,V')$ analogue à $(X,\handwritten{\unsure{U}},G,U')$, et un morphisme $f:X\handwritten{\longrightarrow}Y$ tel que $f(U)\handwritten{\longrightarrow}f(V)$, un morphisme $f':U'\handwritten{\longrightarrow}V'$ compatible avec $U\handwritten{\longrightarrow}V$ induit par $f$, un morphisme de groupes $\handwritten{\phi}:G\handwritten{\longrightarrow}H$ tel que \deleted{$\unsure{f'(gx')=}$} $f'$ soit compatible avec les actions de $G$ et $H$ via $\handwritten{\phi}$, \deleted{\unsure{et}}

Le cas le plus courant sera celui où on se donne un point géométrique $a$ de $U$, \inserted{d'où} un point géométrique $b$ \inserted{$=f(a)$} de $V$, où $U'$ et $V'$ sont respectivement des quotients des revêtements universels de $U$ et $V$ basés en $a$ et $b$, $G$ et $H$ étant respectivement des quotients de $\handwritten{\pi}_{1}(U,a)$ et $\handwritten{\pi}_{1}(V,b)$, l'homomorphisme $f':U'\handwritten{\longrightarrow}V'$ étant induit par les homomorphismes canoniques sur les revêtements universels en $a$ resp. $b$\deleted{\unsure{;} \illegible{2} \unsure{ainsi} \unsure{donc} déterminé de}, \inserted{\deleted{ainsi que $\handwritten{\phi}$}} \deleted{façon unique $G$ et $H$ étant les \illegible{1}} $\handwritten{\phi}:G\handwritten{\longrightarrow}H$ étant induit par $f_{\handwritten{*}}:\handwritten{\pi}_{1}(U,a)\handwritten{\longrightarrow}\handwritten{\pi}_{1}(\unsure{B|V},b)$ ; de cette façon $f'$ et $\handwritten{\phi}$ sont déterminés sans ambiguïté par $f$, et leur existence signifie simplement que $f_{\handwritten{*}}:\handwritten{\pi}_{1}(U,a)\handwritten{\longrightarrow}\handwritten{\pi}_{1}(V,b)$ veut bien passer au quotient en un $\handwritten{\phi}:G\handwritten{\longrightarrow}H$ (i.e. applique un certain sous-groupe dans un certain sous-groupe). Considérons d'autre part le localisé strict $\handwritten{\overline}{Y}$ de $Y$ en le point géométrique $f(\handwritten{\bar}{x})$ de $Y$, d'où \deleted{\illegible{1}} $\handwritten{\bar}{V},\handwritten{\bar}{V}'$ comme dessus; supposons donné un prolongement $\handwritten{\overline}{Y'}$ de $\handwritten{\bar}{V}'$ en un schéma à groupe d'opérateurs $H$ sur $\handwritten{\bar}{Y}$, et supposons donné enfin un homomorphisme $\handwritten{\overline}{f'}:\handwritten{\overline}{X'}\handwritten{\longrightarrow}\handwritten{\overline}{Y'}$ qui prolonge l'homomorphisme \deleted{\unsure{ft|f'}} $\handwritten{\overline}{U'}\handwritten{\longrightarrow}\handwritten{\overline}{V'}$ induit par $f'$. Pour tout point $\handwritten{\unsure{\xi|\zeta}}$ de $\handwritten{\overline}{X'}_{\handwritten{\bar}{x}}$, $\handwritten{\overline}{f'}(\handwritten{\xi})$ est alors un point\deleted{\unsure{u|s}} de $\handwritten{\overline}{Y'}_{\handwritten{\bar}{y}}$, et on peut considérer le groupe d'inertie correspondant $H_{\handwritten{\overline}{f'}(\handwritten{\xi})}$. On a alors la relation évidente
\[
\handwritten{\phi}(G_{\handwritten{\xi}})\handwritten{\subset}H_{\handwritten{\overline}{f'}(\handwritten{\xi})}\ .
\tag*{\handwritten{(2.1)}}
\]
Dans le cas où le\deleted{\unsure{s}} groupe $G$ resp. $H$ opère transitivement sur $\handwritten{\overline}{X'}_{\handwritten{\bar}{x}}$ resp. $\handwritten{\overline}{Y'}_{\handwritten{\overline}{y'}}$, on peut donc dire que $\handwritten{\phi}$ \inserted{applique tout groupe \deleted{de}} \deleted{transforme la classe de conjugaison de sous-groupes} d'inertie de $G$ rel. à $x$ \inserted{\handwritten{\uline{dans}} un groupe d\deleted{e}\handwritten{'}} \deleted{\unsure{en|à} la classe de conj. des sous-groupes} d'inertie de $H$ rel à $y=f(x)$.

\deleted{D'autre part} Il reste alors à expliciter l'homomorphisme $\unsure{\handwritten{\overline}{f'}|\handwritten{\bar f}}:\handwritten{\overline{X'}\longrightarrow\overline{Y'}}$.

% Pencil enclosure around the paragraph beginning ``Le cas'', with a continuation arrow to the last paragraph.
\figure{\handwritten{\begin{tikzpicture}[x=0.5cm,y=0.5cm]
\draw (15,18) .. controls (14.6,18.2) and (14.8,18.5) .. (14,18.4) -- (1,18.4) .. controls (0.6,18.7) and (0.4,18) .. (0.4,17.6) -- (0.3,14) .. controls (0.8,12.5) and (0.5,11.4) .. (0.5,10.5) .. controls (1,10) and (1.3,10.2) .. (2,10.2) -- (7.4,10.2) .. controls (8.1,10.2) and (8.2,10.4) .. (8.2,10.9) -- (14.8,11.1);
\draw[->] (0.5,10.5) .. controls (-0.4,9) and (-0.3,3) .. (0.3,0.7) .. controls (0.5,0.1) and (0.9,0.5) .. (1.5,0.1);
\draw (-1.3,17.4) .. controls (-1.7,17.1) and (-1.7,16.9) .. (-1.5,16.6);
\end{tikzpicture}}}

\scanpage{120}
% archive sheet number: 119

En pratique, il sera vrai que $\handwritten{\overline}{U'}$ est schématiquement dense dans $\handwritten{\overline}{X'}$\inserted{et $\handwritten{\overline}{Y'}$ séparé sur $\handwritten{\bar}{Y}$}, de sorte que $\handwritten{\overline}{f'}$ sera déterminé de façon unique en termes des autres données, comme l'unique prolongement de $\handwritten{\overline}{U'}\handwritten{\longrightarrow}\handwritten{\overline}{V'}$ induit par $f'$; ce sera le cas en particulier dans le cas envisagé plus haut où $\handwritten{\overline}{X'}$ se déduit de \deleted{$\handwritten{\bar}{\unsure{U}}$} $\handwritten{\overline}{U'}$ par normalisation. Bien \unsure{entendur|entendu}, pour pouvoir conclure\inserted{\handwritten{alors}} que tout sous-groupe d'inertie de $G$ rel. à $x$ s'envoie dans un sous-groupe d'inertie de $H$ rel. à $y$, il faudra avoir vérifié l'\uline{existence} d'un \unsure{ppolongement|prolongement} $\handwritten{\overline}{f'}$ de $\handwritten{\overline}{U'}\handwritten{\longrightarrow}\handwritten{\overline}{V'}$\inserted{\handwritten{(la commutation aux actions de $G$ \unsure{dev.|sera} alors \uline{automatique})}}. \handwritten{cela \unsure{marchera}, en interprétant $f'$ comme $\overline{X'}\longrightarrow\overline{Y'}\mathbin{\times}_{\bar Y}\overline{X'}$, lorsque $\overline{Y'}$ est \uline{entier} \deleted{\illegible{1}} sur $\bar Y$, cf. EGA II 6.1.14.}

\handwritten{3.} Revenons à la situation de 1), supposons que $\underline{O}_{X,x}$ soit un anneau local régulier, et que $\handwritten{\bar}{U}$ soit le complémentaire dans le schéma stricte\deleted{\unsure{ment}}ment local régulier $\handwritten{\bar}{X}$ d'un diviseur \inserted{$D$} à croisements normaux. On est en particulier dans le cas favorable où la définition de $\handwritten{\overline}{X'}$ est immédiate par normalisation. Supposons que la ramification de $\handwritten{\overline}{U'}$ en les points maximaux de $\handwritten{\bar}{D}$ soit \uline{modérée}. On sait alors, par le calcul local bien connu, que pour tout $\handwritten{\unsure{\xi|\zeta}\in}\handwritten{\overline}{X'}_{\handwritten{\bar}{x}}$, le groupe d'inertie correspondant $G$ s'identifie canoniquement \deleted{au groupe} à un groupe quotient du groupe
\margin{\handwritten{$C=$ l'ensemble des composantes irred. de $\bar D$}}
\[
I=\handwritten{\prod}_{\handwritten{\unsure{\ell|l}\ne}p}\underline{Z}_{\handwritten{\ell}}(1)(k(\handwritten{\bar}{x}))^C=\handwritten{\prod}_{\handwritten{\ell\ne}p}T_{\handwritten{\ell}}(k(\handwritten{\bar}{x}\unsure{)}\ )^C\ ,
\tag*{(3.1)}
\]
$p$ étant la caractéristique résiduelle en $x$. De p\unsure{l|m}us, ce même calcul local montre que le facteur de ce groupe d'inertie correspondant à la composante $\handwritten{\bar}{D}_c$ $(c\handwritten{\in}C)$ du diviseur de ramification $\handwritten{\bar}{D}$ coïncide aussi avec le groupe d'inertie relatif à n'importe quel point de $\handwritten{\overline}{X'}$ au dessus du point maximal de $\handwritten{\bar}{D}_c$ (ou de tout autre point de $\handwritten{\bar}{D}_c$ qui n'appartient à aucune autre composante irréductible de $\handwritten{\bar}{D}$) qui générise $\handwritten{\xi}$ ; comme il y a toujours de tels points, cela montre que la classe de conjugaison du facteur en question peut aussi s'interpréter comme la \deleted{\unsure{grou}} classe des groupes d'inertie en les points envisagés de $\handwritten{\bar}{D}_c$ ; en particulier, cette classe de conjugaison ne dépend que de la composante irréductible de $D=X-U$ \deleted{\unsure{,} image} de point générique l'image du point générique de $\handwritten{\bar}{D}_{\handwritten{c}}$.

\scanpage{121}
% archive sheet number: 120

\handwritten{4.} Revenons à la situation de 2, et supposons que la structure de $(\handwritten{\bar}{X},\handwritten{\bar}{U})$ soit celle envisagée dans 3, et que de même $\handwritten{\bar}{Y}$ soit un schéma local régulier et $\handwritten{\bar}{V}=\handwritten{\bar}{Y}-\operatorname{supp}\handwritten{\bar}{E}$, $\handwritten{\bar}{E}$ étant un diviseur à croisements normaux sur $\handwritten{\bar}{Y}$. Soit $C'$ l'ensemble des composantes irréductibles de $\handwritten{\bar}{E}$, On voit alors que pour tout $c'\handwritten{\in}C'$, l'image inverse de $\handwritten{\bar}{E}_{c'}$ est un diviseur sur $\handwritten{\bar}{X}$ dont le support est contenu dans celui de $\deleted{\unsure{E}}\handwritten{\bar}{D}$, donc est de la forme
\[
\handwritten{\bar}{f}(\handwritten{\bar}{E}_{c'})=\handwritten{\sum}_{c}n_{c,c'}\handwritten{\bar}{D}_c\ ,
\]
où pour $c'$ fixé, les $n_{c,c'}$ sont des entiers $\handwritten{\geq}0$ non tous nuls. Ces notations fixées, un calcul immédiat \deleted{\unsure{direct}} d'images inverses de revêtements kummériens montre que pour deux points correspondants $\handwritten{\unsure{\xi|\zeta}}$ et $\handwritten{\unsure{\eta|y}}$, \deleted{\unsure{l'homom}} on a commutativité dans le diagramme
\diagram{\(\handwritten{(4.1)}\qquad\)
\begin{tikzcd}[ampersand replacement=\&,column sep=large,row sep=large]
\&\handwritten{\deleted{\unsure{H_\xi|G_\xi}}}\&G_{\handwritten{\xi}}^{\handwritten{\deleted{\unsure{\xi|x}}}}\handwritten{\arrow}[r,"{\handwritten{\phi}}","{\deleted{\unsure{N|\phi}}}"']\&H_{\handwritten{\eta}}\\
\handwritten{\unsure{-}}\&{}\handwritten{\arrow}[u,gray]\&\handwritten{\prod}_{\handwritten{\unsure{\ell|l}}}\underline{Z}_{\handwritten{\ell}}(1)(\handwritten{\bar}{x})^C\handwritten{\arrow}[u]\handwritten{\arrow}[r,"N"]\&\handwritten{\prod}^{\handwritten{\unsure{\prime|\prime\prime}}}_{\handwritten{\ell}}\underline{Z}_{\handwritten{\ell}}(1)(\handwritten{\bar}{y})^{C'}\handwritten{\arrow}[u]
\end{tikzcd}\(\ ,\)}
les flèches \deleted{horizontales}\inserted{verticales} étant les épimorphismes canoniques, et $N$ étant l'homomorphisme défini par la matrice $N=(n_{c,c'}\handwritten{)}_{(c,c')\handwritten{\in}C\mathbin{\times}C'}$, compte tenu que l'homomorphisme $k(\handwritten{\bar}{y})\handwritten{\longrightarrow}k(\handwritten{\bar}{x})$ induit des isomorphismes
\[
\underline{Z}_{\handwritten{\ell}}(1)(\handwritten{\bar}{y})\handwritten{\xrightarrow{\sim}}\underline{Z}_{\handwritten{\ell}}(1)(\handwritten{\bar}{x})\ .
\]
Pour que \inserted{l'homomorphisme} $N$ soit un isomorphisme, il est donc nécessaire et suffisant que \inserted{$\operatorname{card}C=\operatorname{card}C'$ et que} \deleted{\unsure{dét N}} l'entier ordinaire $\text{dét }N$ soit égal à $\pm p^r$, où $p$ est l'exposant caractéristique de $k(x)$ et $k(y)$; dans ce cas, on conclut donc que les \inserted{sous-}groupes d'inertie de $H$ en $y$ sont les conjugués des images des \inserted{sous-}groupes d'inertie de $G$ en $x$.

\margin{\handwritten{$\unsure{E'|\bar E}\unsure{=|\longrightarrow}\unsure{Y|X}^{\unsure{\vee}}$}}
5. Supposons en particulier que l'on ait $V=Y-\operatorname{supp}E$, $E$ étant un diviseur \inserted{$\handwritten{\geq}0$} sur $Y$, \deleted{tel que} $U=f^{\handwritten{-1}}(V)$, et que \deleted{pour tout $x\handwritten{\in}f^{\handwritten{-1}}(E)$, posant} les conditions suivantes soient satisfaites:

\inserted{\handwritten{a) Pour $x\in f^{-1}(\operatorname{supp}E)$, $y=f(x)$, $\underline{O}_{X,x}$ et $\underline{O}_{Y,y}$ sont noethériens.}}

\deleted{\unsure{a}}\handwritten{b}) Le diviseur \inserted{\handwritten{$D=$}} $f^{\handwritten{*}}(E)$ est défini, c'est un sous-schéma \uline{régulier} de

\scanpage{122}
% archive sheet number: 121

$\mathbb{X}$ \deleted{et}\inserted{\handwritten{(}donc} $X$ est régulier en les points de son support\handwritten{)},\deleted{$U=f^{\handwritten{-1}}(V)=X-\operatorname{supp}D$}\handwritten{.}
\margin{\handwritten{(Condition de transversalité)}}

\handwritten{b}) $\operatorname{supp}E$ est irréductible, et il est régulier en les points de $f(\operatorname{supp}D)=\operatorname{supp}E\handwritten{\cap}f(X)$ (donc $Y$ est régulier en ses points).

Alors pour tout $x\handwritten{\in}\operatorname{supp}D$, \deleted{\unsure{l'image des}}\inserted{les} sous-groupes d'inertie de $G$ en $x$ \unsure{ont|sont} comme images par $\handwritten{\phi}:G\handwritten{\longrightarrow}H$ des sous-groupes de $H$ qui sont des sous-groupes d'inertie de $H$ relativement au point maximal de $\operatorname{supp}E$. Par suite, pour toute représentation linéaire de $G$ (dans un module $M$ sur quelque anneau) \deleted{\unsure{sur}} qui veut bien se factoriser par une représentation linéaire \uline{de} $H$, les images dans $\operatorname{Aut}(M)$ des \inserted{sous-}groupes d'inertie \deleted{\unsure{dans}} de $G$ relatifs aux divers points de $\operatorname{supp}D$ (ou encore, relatifs aux diverses composantes irréductibles de $\operatorname{supp}D$, cela revient au même en vertu de 3\handwritten{)} sont conjugués dans $\operatorname{Aut}(M)$.

\inserted{\uline{Remords\handwritten{\unsure{:}} au n\textsuperscript{o}}\handwritten{s3,4}.} On peut demander \deleted{\unsure{un}}\handwritten{\unsure{des|du}} résultat\handwritten{\unsure{s}} plus précis, en \deleted{\uline{Application}} considérant pour tout $\handwritten{\unsure{\xi|\zeta}}$ comme ci-dessus (au dessus d'un point géométrique\inserted{$\handwritten{\bar}{x}$} de $\operatorname{supp}D$) l'homomorphisme correspondant $\underline{Z}_{\handwritten{\unsure{\ell|l}}}(1)(\handwritten{\bar}{x})\handwritten{\longrightarrow}G$ dont l'image est le sous-groupe d'inertie correspondant, et en se demandant comment cet homomorphisme varie avec le choix de $\handwritten{\xi}$ \uline{et de} $\handwritten{\bar}{x}$. Pour $\handwritten{\xi}$ variable, la variation se fait évidemment modulo composition avec les automorphismes intérieurs de $G$, \deleted{Pour donner un sens à la question} donc pour $\handwritten{\bar}{x}$ fixé, on trouve une classe de conjugaison d'homomorphismes
\[
\underline{Z}_{\handwritten{\ell}}(1)(\handwritten{\bar}{x})\handwritten{\longrightarrow}G\ .
\]
\deleted{pour}\inserted{\handwritten{Étudions}} \deleted{L}\handwritten{l}a variation de cet homomorphisme\handwritten{,} pour $\handwritten{\bar}{x}$ variable dans une composante \deleted{de $D$}\inserted{$D_i$ \handwritten{de}} l'ensemble des points réguliers de $D$\inserted{à car. $\handwritten{\ne\ell}$}; d'après ce qu'on a vu au n\textsuperscript{o}\unsure{3|2}, on passe d'un tel homomorphisme à un autre en composant avec un quelconque des isomorphismes $\underline{Z}_{\handwritten{\ell}}(1)(\handwritten{\bar}{x})\handwritten{\longrightarrow}\underline{\unsure{\mathbb{Z}|E|Z}}_{\handwritten{\ell}}(1)(\handwritten{\overline}{x'})$ associé à une classe de chemins de $\handwritten{\bar}{x}$ à $\handwritten{\overline}{x'}$. Si cet isomorphe ne dépend pas du choix de cette classe de chemins, i.e. si le\deleted{action du groupe fondamental} faisceau \unsure{l|1}-adique de \handwritten{T}ate $\underline{Z}_{\handwritten{\ell}}(1)$ sur $D(\handwritten{\ell})$ est trivial (ce qui se voit déjà sur la

% Contour separating the Remords passage from the preceding paragraph; doubled left boundary.
\figure{\handwritten{\begin{tikzpicture}[x=0.6cm,y=0.6cm]
\draw (15,12.5) .. controls (12,12.2) and (10,12.4) .. (5,12.4) -- (4.7,11.5) -- (3.3,11.5) .. controls (1.8,11.1) and (1.4,11.4) .. (0,11.3) .. controls (-0.2,10) and (0.2,7) .. (-0.1,6) .. controls (-0.3,4) and (0.3,2.3) .. (0.1,2) .. controls (-0.2,1) and (0,0.4) .. (0.2,0);
\draw (0,11.3) .. controls (-0.1,8.5) and (-0.2,8) .. (-0.1,6) .. controls (0,3.8) and (0,1.5) .. (0.25,0.4);
\end{tikzpicture}}}

\scanpage{123}
% archive sheet number: 122

restriction dudit faisceau au \deleted{corps résiduel} point générique de $D_i$, alors on peut identifier entre eux tous les groupes $\underline{Z}_{\handwritten{\unsure{\ell|l}}}(1)(\handwritten{\bar}{x})$ envisagés à un même groupe $\underline{\mathbb{Z}}_{\handwritten{\ell}}(1)(D_i)$ (qu'on peut interpréter comme groupe des sections de $\underline{Z}_{\handwritten{\ell}}(1)$ sur $D_i$ ou sur le point générique de $D_i$), et on trouve une classe \inserted{de conjugaison} bien déterminée d\deleted{\unsure{r|e}}'homomorphismes $\underline{Z}_{\handwritten{\ell}}(1)(D_i)\handwritten{\mathbin{-}}G$. (On trouve un résultat analogue pour les homomorphismes de produits de groupes $\underline{Z}_{\handwritten{\ell}}(1)$ envisagés dans le n\textsuperscript{o}3). Bien entendu, si $\underline{Z}_{\handwritten{\ell}}(1)$ est trivial sur l'ouvert\inserted{$X_{\handwritten{\ell}}$} des points de $X$ \deleted{Lorsque l'on a un $f:X\unsure{\mathbin{-}}Y$ comme au n\textsuperscript{o}\unsure{4|2}} de car. $\handwritten{\ne\ell}$, et si ce dernier est connexe, on peut identifier les groupes $\underline{Z}_{\handwritten{\ell}}(1)(D_i)$ au même groupe \inserted{\handwritten{\unsure{à|2}}} $\underline{Z}_{\handwritten{\ell}}(1)(X_{\handwritten{\ell}})$, \deleted{et on} soit $\underline{Z}_{\handwritten{\ell}}(1)$, et on trouve \deleted{un homomorphisme cano} pour chaque $i$ une\deleted{\unsure{homomo}} classe de conjugaison d'\unsure{homomorphismens|homomorphisnens} $\underline{Z}_{\handwritten{\ell}}(1)\handwritten{\longrightarrow}G$\handwritten{.} En général, celles-ci \deleted{sont} peuvent être distinctes pour des $i$ distincts, bien entendu. Mais relativement à un\deleted{e} morphisme $f:X\handwritten{\longrightarrow}Y$ comme dans les n\textsuperscript{o}2,\handwritten{4}, la formation des classes\inserted{de conjugaison} \deleted{précédentes} d'homomorphismes\inserted{ci-dessus} a un caractère évident de fonctorialité (qu'il y aurait lieu d'énoncer de façon précise, dans le cadre général du n\textsuperscript{o} 4).

% Left and lower contour enclosing the continuation of the Remords passage.
\figure{\handwritten{\begin{tikzpicture}[x=0.65cm,y=0.65cm]
\draw (0,11.5) .. controls (0.2,10.8) and (0,9.6) .. (0.3,9) .. controls (0.4,6) and (0.2,3) .. (0.6,0);
\draw (0.7,-0.05) .. controls (1.8,0.1) and (3,0) .. (5,0) -- (9,0) .. controls (9.4,0) and (9.2,0.4) .. (9.7,0.35) .. controls (11.5,0.1) and (13,0.2) .. (15,0.5);
\end{tikzpicture}}}

Revenant alors aux conditions du présent n\textsuperscript{o}5, on peut préciser de la façon suivante l'assertion faite sur les \deleted{grou} classes de conjugaison de sous-groupes d'inertie en des assertions sur les classes de conjugaisons d'\uline{homomorphismes} de \uline{modules de Tate}: supposons que $Y_{\handwritten{\ell}}$ soit connexe et que $Z_{\handwritten{\ell}}(1)$ soit trivia\unsure{l|b} sur $Y_{\handwritten{\ell}}$, \deleted{\unsure{d'où|d'on} un isomorphisme désignant} de sorte qu'on peut identifier $Z_{\handwritten{\ell}}(1)(Y_{\handwritten{\ell}})$ et les divers $\underline{Z}_{\handwritten{\ell}}(1)(D_i)$ à un même module \handwritten{(}libre de rang 1\handwritten{)}$\underline{Z}_{\handwritten{\ell}}(1)$ sur $\underline{Z}_{\handwritten{\ell}}$. Considérons alors les classes de conjugaison d'homomorphismes composés
\diagram{\begin{tikzcd}[ampersand replacement=\&,column sep=huge]
\underline{Z}_{\handwritten{\ell}}(1)\handwritten{\arrow}[r,"{\text{inertie en }D_i}"]\&G\handwritten{\arrow}[r,"{\handwritten{\phi}}"]\&H
\end{tikzcd}\(.\)}

Ces classes sont conjuguées entre elles.

\scanpage{125}
\mspage{\handwritten{\textcircled{1}}}
% archive sheet number: 124

\handwritten{\unsure{Vallouin|Valloire|Vallouise} 18.\unsure{6|8}.1959}\\

Mon cher Serre;\\

Tate m'a écrit de son côté \deleted{pour} sur ses histoires de courbes elliptiques, et pour me demander si j'avais des idées sur une définition globale des variétés \unsure{anplytiques|analytiques} sur des corps valués complets. Je dois avouer que je n'ai pas du tout compris pourquoi ses résultats suggèreraient l'existence d'une telle définition, et suis encore sceptique. Je n'ai pas non plus l'impression d'avoir rien compris à son théorème, qui \deleted{pour moi} ne fait que exhiber par des formules brutales un certain isomorphisme de groupes analytiques; \deleted{et} on conçoit que d'autres formules tout aussi explicites en donneraient un autre pas plus mauvais (sauf preuve du contraire !)

J'ai réfléchi un petit peu à la partie 'infinitésimale' du groupe fondamental, juste assez pour me convaincre que ça existe et est raisonnable. Voici un \unsure{contexte} (certainement insuffisant d'ailleurs) où ça marche. On a un schéma $S$\inserted{\handwritten{connexe}} (par exemple un schéma algébrique sur un corps $k$), une catégorie auxiliaire $\underline{C}$ (par exemple la catégorie des schémas algébriques finis sur $k$), une catégorie \deleted{de} $\underline{G}$ dont les objets sont des $\underline{C}$-groupes, et les morphismes des morphismes de $\underline{C}$-groupes (par exemple les groupes algébriques finis sur $k$). On \unsure{suppoce|suppose} que dans $\underline{C}$ les produits \deleted{\unsure{finis}}\handwritten{fibrés} existent, et que $\underline{G}$ satisfait aux conditions suivantes: (i) $\underline{\mathbb{G}}$ est stable par produits, et si $u,v:G\handwritten{\longrightarrow}G'$ sont des morphismes dans $\underline{G}$, alors le noyau du couple $(u,v)$ i.e. le sous-groupe \deleted{\illegible{1}} maximum de $G$ où $u,v$ coïncidant -qui existe, car s'exprime par un produit fibré- est dans $\underline{G}$. (ii) Si $u:G\handwritten{\longrightarrow}G'$ est un morphisme dans $\underline{G}$, \deleted{alors il se décompose} le groupe image existe, est isomorphe à un quotient de $G$ comme il se doit, et est dans $\underline{G}$. (iii) Toute suite décroissante de sous-groupes \deleted{\unsure{d'un}} $\handwritten{\in}\underline{G}$ d'un $G\handwritten{\in}\underline{C}$ est stationnaire. On suppose donné de plus un foncteur covariant \deleted{$\unsure{F:G|F:C}$} $F$ de $\underline{G}$ dans les \deleted{\unsure{G-groupes|S-groupes}} schémas en groupes sur $S$ (par exemple $G\handwritten{\longrightarrow}S\mathbin{\times}_{k}G$), et on suppose : (iv) Le foncteur $F$ commute aux produits, \inserted{\handwritten{aux}} noyaux de \unsure{coupbes|couples} de morphismes, images (on pourra dire que $\mathbb{F}$ est ``exact''). (v) Si $H$ est \deleted{un} de la forme $F(G)$ $(G\handwritten{\in}\underline{G})$ alors $H$\inserted{\handwritten{est ``spécial'', i.e.}} admet une suite exacte
